\chapter{Formules et équations trigonométriques}
\textit{Tu connais déjà le cercle trigonométrique, les valeurs remarquables et les
fonctions $\cos$, $\sin$, $\tan$. Ce chapitre te donne les \emph{formules} qui
relient ces fonctions entre elles — addition, duplication, linéarisation — puis te
montre comment résoudre les \emph{équations} et \emph{inéquations} trigonométriques.
Ce sont des outils incontournables au Probatoire comme dans toute la suite.}

\section{Rappels et conventions}

\begin{aretenir}[Ce qu'il faut avoir en tête]
Un angle est mesuré en \textbf{radians}. Sur le cercle trigonométrique (de rayon
$1$), un réel $x$ repère un point $M$ tel que $\cos x$ est son abscisse et $\sin x$
son ordonnée. On rappelle l'\textbf{identité fondamentale}
\[ \cos^2 x+\sin^2 x=1, \]
ainsi que $\tan x=\dfrac{\sin x}{\cos x}$, définie pour $x\neq \dfrac{\pi}{2}+k\pi$
($k\in\Z$). Les fonctions $\cos$ et $\sin$ sont $2\pi$-périodiques, $\tan$ est
$\pi$-périodique.
\end{aretenir}

\begin{illus}
\begin{tikzpicture}[scale=2.1]
  \draw[->,onbline] (-1.35,0)--(1.4,0) node[right,onbmuted]{$x$};
  \draw[->,onbline] (0,-1.35)--(0,1.4) node[above,onbmuted]{$y$};
  \draw[thick,onbmuted] (0,0) circle (1);
  \coordinate (O) at (0,0);
  \coordinate (A) at (1,0);
  \coordinate (M) at (50:1);
  \draw[onbo,very thick] (O)--(M);
  \filldraw[onbo] (M) circle (0.6pt) node[above right,onbink]{$M$};
  \draw[dashed,onbmuted] (M)--(0.643,0) node[below,onbink]{$\cos x$};
  \draw[dashed,onbmuted] (M)--(0,0.766) node[left,onbink]{$\sin x$};
  \draw pic[draw=onbgreen,angle radius=0.4cm,"$x$"{onbgreen}] {angle=A--O--M};
\end{tikzpicture}
\fig{Figure — sur le cercle de rayon $1$, $\cos x$ est l'abscisse et $\sin x$ l'ordonnée du point $M$ repéré par $x$.}
\end{illus}

\section{Formules d'addition}

\begin{theoreme}[Formules d'addition]
Pour tous réels $a$ et $b$ :
\begin{align*}
  \cos(a+b) &= \cos a\cos b-\sin a\sin b, &
  \cos(a-b) &= \cos a\cos b+\sin a\sin b,\\
  \sin(a+b) &= \sin a\cos b+\cos a\sin b, &
  \sin(a-b) &= \sin a\cos b-\cos a\sin b.
\end{align*}
Pour la tangente (lorsque tous les termes sont définis) :
\[
  \tan(a+b)=\frac{\tan a+\tan b}{1-\tan a\tan b},\qquad
  \tan(a-b)=\frac{\tan a-\tan b}{1+\tan a\tan b}.
\]
\end{theoreme}

\begin{remarque}
Pour mémoriser : pour $\cos$, on garde le \emph{même} type de produit
($\cos\cos$, $\sin\sin$) avec le signe \emph{opposé} à celui de l'angle ; pour
$\sin$, on \emph{croise} ($\sin\cos$, $\cos\sin$) avec le \emph{même} signe.
\end{remarque}

\begin{exemple}
Calculons $\cos\dfrac{\pi}{12}$. Comme $\dfrac{\pi}{12}=\dfrac{\pi}{3}-\dfrac{\pi}{4}$ :
\[
  \cos\frac{\pi}{12}=\cos\frac{\pi}{3}\cos\frac{\pi}{4}+\sin\frac{\pi}{3}\sin\frac{\pi}{4}
  =\frac12\cdot\frac{\sqrt2}{2}+\frac{\sqrt3}{2}\cdot\frac{\sqrt2}{2}
  =\frac{\sqrt2+\sqrt6}{4}.
\]
\end{exemple}

\begin{propriete}[Angles associés (cas particuliers utiles)]
En prenant $b=\dfrac{\pi}{2}$ ou $b=\pi$ dans les formules d'addition, on retrouve :
\[
  \cos\!\Big(\frac{\pi}{2}-x\Big)=\sin x,\quad
  \sin\!\Big(\frac{\pi}{2}-x\Big)=\cos x,\quad
  \cos(\pi-x)=-\cos x,\quad
  \sin(\pi-x)=\sin x.
\]
\end{propriete}

\section{Formules de duplication}

\begin{theoreme}[Formules de duplication]
Pour tout réel $a$ :
\[
  \cos 2a=\cos^2 a-\sin^2 a,\qquad
  \sin 2a=2\sin a\cos a,\qquad
  \tan 2a=\frac{2\tan a}{1-\tan^2 a}.
\]
En utilisant $\cos^2 a+\sin^2 a=1$, la formule de $\cos 2a$ s'écrit aussi
\[
  \cos 2a=2\cos^2 a-1=1-2\sin^2 a.
\]
\end{theoreme}

\begin{exemple}[Démonstration à partir de l'addition]
Il suffit de poser $b=a$ dans les formules d'addition. Pour le cosinus :
\[
  \cos 2a=\cos(a+a)=\cos a\cos a-\sin a\sin a=\cos^2 a-\sin^2 a.
\]
Pour le sinus :
\[
  \sin 2a=\sin(a+a)=\sin a\cos a+\cos a\sin a=2\sin a\cos a.
\]
De même $\tan 2a=\tan(a+a)=\dfrac{\tan a+\tan a}{1-\tan a\tan a}=\dfrac{2\tan a}{1-\tan^2 a}$.
\end{exemple}

\begin{exemple}
Connaissant $\cos\dfrac{\pi}{6}=\dfrac{\sqrt3}{2}$, retrouvons $\cos\dfrac{\pi}{3}$ :
\[
  \cos\frac{\pi}{3}=\cos\!\Big(2\cdot\frac{\pi}{6}\Big)=2\cos^2\frac{\pi}{6}-1
  =2\cdot\frac34-1=\frac12.
\]
\end{exemple}

\section{Formules de linéarisation}

\begin{propriete}[Linéarisation de $\cos^2$ et $\sin^2$]
En isolant $\cos^2 a$ puis $\sin^2 a$ dans les deux écritures de $\cos 2a$, on
obtient :
\[
  \cos^2 a=\frac{1+\cos 2a}{2},\qquad
  \sin^2 a=\frac{1-\cos 2a}{2}.
\]
\end{propriete}

\begin{remarque}
\emph{Linéariser}, c'est remplacer une puissance ($\cos^2 a$, $\sin^2 a$) par une
expression où les fonctions trigonométriques n'apparaissent qu'à la puissance $1$
(ici $\cos 2a$). C'est essentiel pour calculer certaines intégrales en Terminale.
\end{remarque}

\begin{exemple}
Calculons $\sin^2\dfrac{\pi}{8}$. Avec $a=\dfrac{\pi}{8}$, donc $2a=\dfrac{\pi}{4}$ :
\[
  \sin^2\frac{\pi}{8}=\frac{1-\cos\frac{\pi}{4}}{2}
  =\frac{1-\frac{\sqrt2}{2}}{2}=\frac{2-\sqrt2}{4}.
\]
\end{exemple}

\section{Transformation de $a\cos x+b\sin x$}

\begin{theoreme}[Forme $R\cos(x-\varphi)$]
Soit $a$ et $b$ deux réels non tous deux nuls. En posant $R=\sqrt{a^2+b^2}$, il
existe un réel $\varphi$ tel que, pour tout réel $x$,
\[
  a\cos x+b\sin x=R\cos(x-\varphi),
\]
où $\varphi$ est déterminé par $\cos\varphi=\dfrac{a}{R}$ et $\sin\varphi=\dfrac{b}{R}$.
\end{theoreme}

\begin{exemple}
En développant $R\cos(x-\varphi)=R(\cos\varphi\cos x+\sin\varphi\sin x)$ et en
identifiant à $a\cos x+b\sin x$, on retrouve $R\cos\varphi=a$ et $R\sin\varphi=b$ ;
d'où $R^2(\cos^2\varphi+\sin^2\varphi)=a^2+b^2$, soit $R=\sqrt{a^2+b^2}$.
\end{exemple}

\begin{exemple}
Transformons $\cos x+\sqrt3\,\sin x$. Ici $a=1$, $b=\sqrt3$, donc
$R=\sqrt{1+3}=2$. On cherche $\varphi$ avec $\cos\varphi=\dfrac12$ et
$\sin\varphi=\dfrac{\sqrt3}{2}$ : c'est $\varphi=\dfrac{\pi}{3}$. Ainsi
\[
  \cos x+\sqrt3\,\sin x=2\cos\!\Big(x-\frac{\pi}{3}\Big).
\]
\end{exemple}

\begin{aretenir}
Cette transformation est précieuse : elle ramène une somme $a\cos x+b\sin x$ à
\emph{une seule} fonction trigonométrique. On en déduit immédiatement que son
maximum est $R=\sqrt{a^2+b^2}$ et son minimum $-R$, et l'on sait résoudre les
équations $a\cos x+b\sin x=c$.
\end{aretenir}

\section{Équations trigonométriques}

\subsection{Équation $\cos x=a$}

\begin{theoreme}[Résolution de $\cos x=\cos\alpha$]
L'équation $\cos x=a$ admet des solutions si et seulement si $-1\le a\le 1$. En
écrivant $a=\cos\alpha$, on a, pour $k\in\Z$ :
\[
  \cos x=\cos\alpha \iff x=\alpha+2k\pi \ \text{ ou }\ x=-\alpha+2k\pi.
\]
\end{theoreme}

\begin{illus}
\begin{tikzpicture}[scale=2.1]
  \draw[->,onbline] (-1.35,0)--(1.4,0) node[right,onbmuted]{$x$};
  \draw[->,onbline] (0,-1.35)--(0,1.4) node[above,onbmuted]{$y$};
  \draw[thick,onbmuted] (0,0) circle (1);
  \draw[dashed,onbo] (0.5,-1.2)--(0.5,1.2);
  \node[onbo,below right] at (0.5,-1.2){$x=a$};
  \filldraw[onbo] (60:1) circle (0.7pt) node[above right,onbink]{$\alpha$};
  \filldraw[onbo] (-60:1) circle (0.7pt) node[below right,onbink]{$-\alpha$};
  \draw[onbgreen,very thick] (0,0)--(60:1);
  \draw[onbgreen,very thick] (0,0)--(-60:1);
\end{tikzpicture}
\fig{Figure — $\cos x=a$ : les solutions sont les points du cercle d'abscisse $a$, symétriques par rapport à l'axe $(Ox)$.}
\end{illus}

\begin{exemple}
Résolvons $\cos x=\dfrac12$ dans $\R$. Comme $\dfrac12=\cos\dfrac{\pi}{3}$ :
\[
  x=\frac{\pi}{3}+2k\pi \quad\text{ou}\quad x=-\frac{\pi}{3}+2k\pi,\qquad k\in\Z.
\]
Sur $]-\pi\,;\,\pi]$, les solutions sont $\dfrac{\pi}{3}$ et $-\dfrac{\pi}{3}$.
\end{exemple}

\subsection{Équation $\sin x=a$}

\begin{theoreme}[Résolution de $\sin x=\sin\alpha$]
L'équation $\sin x=a$ admet des solutions si et seulement si $-1\le a\le 1$. En
écrivant $a=\sin\alpha$, on a, pour $k\in\Z$ :
\[
  \sin x=\sin\alpha \iff x=\alpha+2k\pi \ \text{ ou }\ x=\pi-\alpha+2k\pi.
\]
\end{theoreme}

\begin{exemple}
Résolvons $\sin x=\dfrac{\sqrt2}{2}$ dans $\R$. Comme $\dfrac{\sqrt2}{2}=\sin\dfrac{\pi}{4}$ :
\[
  x=\frac{\pi}{4}+2k\pi \quad\text{ou}\quad x=\pi-\frac{\pi}{4}+2k\pi=\frac{3\pi}{4}+2k\pi,\qquad k\in\Z.
\]
\end{exemple}

\subsection{Équation $\tan x=a$}

\begin{theoreme}[Résolution de $\tan x=\tan\alpha$]
Pour tout réel $a$, l'équation $\tan x=a$ admet des solutions. En écrivant
$a=\tan\alpha$, on a, pour $k\in\Z$ :
\[
  \tan x=\tan\alpha \iff x=\alpha+k\pi.
\]
\end{theoreme}

\begin{remarque}
La période de $\tan$ étant $\pi$ (et non $2\pi$), les solutions sont espacées de
$\pi$ : il n'y a qu'\emph{une} famille de solutions, contre deux pour $\cos$ et
$\sin$.
\end{remarque}

\begin{exemple}
Résolvons $\tan x=\sqrt3$ dans $\R$. Comme $\sqrt3=\tan\dfrac{\pi}{3}$ :
\[
  x=\frac{\pi}{3}+k\pi,\qquad k\in\Z.
\]
Sur $]-\dfrac{\pi}{2}\,;\,\dfrac{\pi}{2}[$, la seule solution est $\dfrac{\pi}{3}$.
\end{exemple}

\subsection{Résolution sur un intervalle donné}

\begin{aretenir}[Garder les solutions d'un intervalle]
Pour résoudre sur un intervalle $I$ (par exemple $[0\,;\,2\pi[$), on écrit d'abord
les solutions générales en $k$, puis on \emph{teste} les valeurs de $k$
($\dots,-1,0,1,\dots$) et on conserve uniquement celles qui tombent dans $I$.
\end{aretenir}

\begin{exemple}
Résolvons $\cos x=\dfrac12$ sur $[0\,;\,2\pi[$. Les solutions générales sont
$x=\dfrac{\pi}{3}+2k\pi$ et $x=-\dfrac{\pi}{3}+2k\pi$. Dans $[0\,;\,2\pi[$ :
$\dfrac{\pi}{3}$ (pour $k=0$) et $-\dfrac{\pi}{3}+2\pi=\dfrac{5\pi}{3}$ (pour $k=1$).
D'où $S=\Big\{\dfrac{\pi}{3}\,;\,\dfrac{5\pi}{3}\Big\}$.
\end{exemple}

\section{Inéquations trigonométriques simples}

\begin{aretenir}[Lecture sur le cercle]
Pour résoudre une inéquation du type $\cos x\le a$ ou $\sin x\ge a$ sur un
intervalle, on place sur le cercle les solutions de l'\emph{équation}
correspondante, puis on lit \emph{l'arc} qui satisfait l'inégalité :
\begin{itemize}
  \item $\cos x$ se lit sur l'axe horizontal ($\cos x\ge a$ : arc à \emph{droite}
        de la verticale $x=a$) ;
  \item $\sin x$ se lit sur l'axe vertical ($\sin x\ge a$ : arc \emph{au-dessus}
        de l'horizontale $y=a$).
\end{itemize}
\end{aretenir}

\begin{illus}
\begin{tikzpicture}[scale=2.1]
  \draw[->,onbline] (-1.35,0)--(1.4,0) node[right,onbmuted]{$x$};
  \draw[->,onbline] (0,-1.35)--(0,1.4) node[above,onbmuted]{$y$};
  \draw[thick,onbmuted] (0,0) circle (1);
  \draw[dashed,onbo] (0.5,-1.2)--(0.5,1.2);
  \draw[onbo,very thick] (-60:1) arc (-60:60:1);
  \filldraw[onbink] (60:1) circle (0.7pt) node[above right,onbink]{$\frac{\pi}{3}$};
  \filldraw[onbink] (-60:1) circle (0.7pt) node[below right,onbink]{$-\frac{\pi}{3}$};
  \node[onbo] at (1.05,0){};
  \node[onbo,right] at (0.92,0.0){};
\end{tikzpicture}
\fig{Figure — $\cos x\ge\frac12$ : l'arc en gras (abscisses $\ge\frac12$) correspond à $x\in\big[-\frac{\pi}{3}\,;\,\frac{\pi}{3}\big]$ sur $]-\pi\,;\,\pi]$.}
\end{illus}

\begin{exemple}
Résolvons $\cos x\ge\dfrac12$ sur $]-\pi\,;\,\pi]$. L'équation $\cos x=\dfrac12$
donne $x=\pm\dfrac{\pi}{3}$. Le cosinus est l'abscisse du point : il est
$\ge\dfrac12$ sur l'arc situé \emph{à droite} de la verticale, c'est-à-dire pour
\[
  x\in\Big[-\frac{\pi}{3}\,;\,\frac{\pi}{3}\Big].
\]
\end{exemple}

\begin{exemple}
Résolvons $\sin x>\dfrac{\sqrt2}{2}$ sur $[0\,;\,2\pi[$. L'équation $\sin x=\dfrac{\sqrt2}{2}$
donne $x=\dfrac{\pi}{4}$ et $x=\dfrac{3\pi}{4}$. Le sinus est l'ordonnée : il est
\emph{strictement} supérieur à $\dfrac{\sqrt2}{2}$ sur l'arc au-dessus de
l'horizontale, soit
\[
  x\in\Big]\frac{\pi}{4}\,;\,\frac{3\pi}{4}\Big[.
\]
\end{exemple}

\section{Allure des fonctions $\cos$ et $\sin$}

\begin{illus}
\begin{tikzpicture}
\begin{axis}[
  width=12cm,height=5.4cm,
  axis lines=middle,
  xlabel={$x$},ylabel={$y$},
  xlabel style={right},ylabel style={above},
  xmin=-0.4,xmax=6.9,ymin=-1.5,ymax=1.6,
  xtick={1.5708,3.1416,4.7124,6.2832},
  xticklabels={$\frac{\pi}{2}$,$\pi$,$\frac{3\pi}{2}$,$2\pi$},
  ytick={-1,1},
  tick style={onbmuted},
  every axis x label/.style={at={(ticklabel* cs:1)},anchor=west},
  every axis y label/.style={at={(ticklabel* cs:1)},anchor=south},
  clip=true]
  \addplot[onbo,very thick,smooth,samples=120,domain=0:6.2832]{cos(deg(x))};
  \addplot[onbgreen,very thick,smooth,samples=120,domain=0:6.2832]{sin(deg(x))};
  \node[onbo] at (axis cs:0.55,1.25){$\cos$};
  \node[onbgreen] at (axis cs:1.05,1.25){$\sin$};
\end{axis}
\end{tikzpicture}
\fig{Figure — sur $[0\,;\,2\pi]$ : la sinusoïde $\cos$ (orange) et $\sin$ (vert), de période $2\pi$, à valeurs dans $[-1\,;\,1]$.}
\end{illus}

\begin{remarque}
Les courbes de $\cos$ et $\sin$ se déduisent l'une de l'autre par une translation
horizontale de $\dfrac{\pi}{2}$, ce qui traduit l'égalité
$\sin x=\cos\!\big(x-\dfrac{\pi}{2}\big)$.
\end{remarque}

% ════════════════════════════════════════════════════════════
\section{L'essentiel à retenir}

\begin{synthese}
\textbf{Addition.}
\[
  \cos(a\pm b)=\cos a\cos b\mp\sin a\sin b,\qquad
  \sin(a\pm b)=\sin a\cos b\pm\cos a\sin b,
\]
\[
  \tan(a\pm b)=\frac{\tan a\pm\tan b}{1\mp\tan a\tan b}.
\]

\smallskip
\textbf{Duplication.} $\cos 2a=\cos^2 a-\sin^2 a=2\cos^2 a-1=1-2\sin^2 a$ ;\quad
$\sin 2a=2\sin a\cos a$ ;\quad $\tan 2a=\dfrac{2\tan a}{1-\tan^2 a}$.

\smallskip
\textbf{Linéarisation.} $\cos^2 a=\dfrac{1+\cos 2a}{2}$ ;\quad
$\sin^2 a=\dfrac{1-\cos 2a}{2}$.

\smallskip
\textbf{Forme $R\cos(x-\varphi)$.} $a\cos x+b\sin x=R\cos(x-\varphi)$ avec
$R=\sqrt{a^2+b^2}$, $\cos\varphi=\dfrac aR$, $\sin\varphi=\dfrac bR$
(max $R$, min $-R$).

\smallskip
\textbf{Équations} ($k\in\Z$) :
\[
  \cos x=\cos\alpha \iff x=\pm\alpha+2k\pi;\qquad
  \sin x=\sin\alpha \iff x=\alpha+2k\pi \text{ ou } x=\pi-\alpha+2k\pi;
\]
\[
  \tan x=\tan\alpha \iff x=\alpha+k\pi.
\]
\textbf{Méthode} : écrire $a$ sous la forme $\cos\alpha$ / $\sin\alpha$ /
$\tan\alpha$, donner la solution générale, puis garder celles de l'intervalle.

\smallskip
\textbf{Inéquations} : se lisent sur le cercle ($\cos$ en abscisse, $\sin$ en
ordonnée) — on repère l'arc qui convient.
\end{synthese}

% ════════════════════════════════════════════════════════════
\section{Méthodes \& savoir-faire}

\methode{Calculer une valeur exacte avec les formules d'addition}
Décomposer l'angle en somme ou différence d'angles \emph{remarquables}
($\dfrac{\pi}{6},\dfrac{\pi}{4},\dfrac{\pi}{3}$), puis appliquer la formule
correspondante.
\begin{exemple}
$\sin\dfrac{7\pi}{12}=\sin\!\Big(\dfrac{\pi}{3}+\dfrac{\pi}{4}\Big)
=\sin\dfrac{\pi}{3}\cos\dfrac{\pi}{4}+\cos\dfrac{\pi}{3}\sin\dfrac{\pi}{4}
=\dfrac{\sqrt6+\sqrt2}{4}.$
\end{exemple}

\methode{Exprimer $\cos 2a$, $\sin 2a$ à partir d'une donnée}
Connaissant $\sin a$ (ou $\cos a$), utiliser $\cos^2 a+\sin^2 a=1$ pour obtenir
l'autre, puis appliquer les formules de duplication.
\begin{exemple}
Si $\cos a=\dfrac35$ avec $a\in\big]0\,;\,\dfrac{\pi}{2}\big[$, alors
$\sin a=\dfrac45$ (positif), donc $\sin 2a=2\cdot\dfrac45\cdot\dfrac35=\dfrac{24}{25}$
et $\cos 2a=2\Big(\dfrac35\Big)^2-1=-\dfrac{7}{25}$.
\end{exemple}

\methode{Linéariser une expression}
Remplacer $\cos^2 a$ et $\sin^2 a$ par $\dfrac{1\pm\cos 2a}{2}$ pour abaisser le
degré.
\begin{exemple}
$\cos^2 x-\sin^2 x=\dfrac{1+\cos 2x}{2}-\dfrac{1-\cos 2x}{2}=\cos 2x$
(on retrouve bien la duplication).
\end{exemple}

\methode{Transformer $a\cos x+b\sin x$ en $R\cos(x-\varphi)$}
Calculer $R=\sqrt{a^2+b^2}$, puis identifier $\varphi$ par $\cos\varphi=\dfrac aR$
et $\sin\varphi=\dfrac bR$ (un angle remarquable en pratique).
\begin{exemple}
$\sqrt3\,\cos x+\sin x$ : $R=\sqrt{3+1}=2$, $\cos\varphi=\dfrac{\sqrt3}{2}$,
$\sin\varphi=\dfrac12$, d'où $\varphi=\dfrac{\pi}{6}$ et
$\sqrt3\,\cos x+\sin x=2\cos\!\big(x-\dfrac{\pi}{6}\big)$.
\end{exemple}

\methode{Résoudre une équation trigonométrique sur un intervalle}
Ramener à $\cos x=\cos\alpha$, $\sin x=\sin\alpha$ ou $\tan x=\tan\alpha$, écrire la
solution générale en $k$, puis tester les valeurs de $k$ et garder celles de
l'intervalle demandé.
\begin{exemple}
$2\sin x-1=0$ sur $[0\,;\,2\pi[$ : $\sin x=\dfrac12=\sin\dfrac{\pi}{6}$, donc
$x=\dfrac{\pi}{6}+2k\pi$ ou $x=\dfrac{5\pi}{6}+2k\pi$ ;
$S=\Big\{\dfrac{\pi}{6}\,;\,\dfrac{5\pi}{6}\Big\}$.
\end{exemple}

\methode{Résoudre une inéquation par lecture du cercle}
Résoudre l'équation associée pour repérer les « bornes » sur le cercle, puis lire
l'arc qui vérifie l'inégalité ($\cos$ en abscisse, $\sin$ en ordonnée).
\begin{exemple}
$\tan x\le 1$ sur $\big]-\dfrac{\pi}{2}\,;\,\dfrac{\pi}{2}\big[$ : $\tan x=1$ donne
$x=\dfrac{\pi}{4}$ ; comme $\tan$ est croissante sur cet intervalle,
$\tan x\le 1$ équivaut à $x\in\big]-\dfrac{\pi}{2}\,;\,\dfrac{\pi}{4}\big]$.
\end{exemple}
